% ============================================================
% 开题报告专用格式设置
% 正文主体沿用 nwafuthesis 的章、节、正文、图表与页眉页脚体系。
% ============================================================

\graphicspath{{figs/}}

% 学位层次显示文字：由 settings/degree.tex 的 doctor/master 自动决定
\newcommand{\ReportDegreeName}{\ifReportDoctor 博士\else 硕士\fi}

% 占位提示：交稿前建议删除或替换。
\newcommand{\placeholder}[1]{%
  {\color{gray}\small\emph{[待填写：#1]}}%
}

% 固定宽度的封面填写线
\newcommand{\coverline}[2][135pt]{%
  \makebox[#1][c]{\raisebox{0.15ex}{\rule{#1}{0.35pt}}}%
  \hspace{-#1}\makebox[#1][c]{#2}%
}

% 图示占位框
\newcommand{\figureplaceholder}[2][0.88\textwidth]{%
  \begin{figure}[htbp]
    \centering
    \fbox{\parbox[c][5.0cm][c]{#1}{\centering\color{gray}#2}}
    \caption{#2}
  \end{figure}
}

% 正文页眉：根据硕/博选择自动切换。
\newcommand{\setreportheader}{%
  \nwafuhead[C]{%
    \small
    \ifodd\value{page}
      \nouppercase{\leftmark}%
    \else
      西北农林科技大学\ReportDegreeName 研究生开题报告%
    \fi
  }%
  \nwafufoot[C]{\small\thepage}%
}


% --- 封面：字体/字号取自原 Word，坐标取自用户提供的 Word 导出 PDF ---
% 原始 PDF: A4 595.32 x 841.92 pt；以下坐标均采用 PDF 左上角为原点的 pt 值。
% 字体与字号（来自 Word OOXML）：
%   附件1：等线 Light，14 pt（Word 中为粗体）
%   校名/学位论文/开题报告：宋体，18 pt
%   题目标签：黑体，22 pt
%   下部字段与日期：宋体，15 pt
% 英文/数字：Times New Roman。
%
% Windows + XeLaTeX 可直接使用上述系统字体。非 Windows 环境自动回退，
% 但若追求 1:1，请在安装有 SimSun/SimHei/DengXian/Times New Roman 的环境编译。

% --- 封面字体 ---
\IfFileExists{settings/cover-fonts-local.tex}{%
  % 仅用于本地视觉核对，不随模板分发字体文件。
  \input{settings/cover-fonts-local.tex}%
}{%
  \IfFontExistsTF{SimSun}{\newCJKfontfamily\CoverSong{SimSun}}{\newCJKfontfamily\CoverSong{Noto Serif CJK SC}}
  \IfFontExistsTF{SimHei}{\newCJKfontfamily\CoverHei{SimHei}}{\newCJKfontfamily\CoverHei{Noto Sans CJK SC}}
  \IfFontExistsTF{DengXian Light}{\newCJKfontfamily\CoverDeng{DengXian Light}}{%
    \IfFontExistsTF{DengXian}{\newCJKfontfamily\CoverDeng{DengXian}}{\newCJKfontfamily\CoverDeng{Noto Sans CJK SC}}%
  }
  \IfFontExistsTF{DengXian Light}{\newfontfamily\CoverDengAll{DengXian Light}}{%
    \IfFontExistsTF{DengXian}{\newfontfamily\CoverDengAll{DengXian}}{\newfontfamily\CoverDengAll{Noto Sans CJK SC}}%
  }
  \IfFontExistsTF{Times New Roman}{\newfontfamily\CoverTNR{Times New Roman}}{\newfontfamily\CoverTNR{Liberation Serif}}
}

% PDF 坐标 -> TikZ current page 坐标的辅助宏。
\newcommand{\pdfbasewest}[4]{%
  \node[anchor=base west,inner sep=0pt,outer sep=0pt] at
    ([xshift=#1bp,yshift=-#2bp]current page.north west) {#4};%
}
\newcommand{\pdfbasecenter}[4]{%
  \node[anchor=base,inner sep=0pt,outer sep=0pt] at
    ([xshift=#1bp,yshift=-#2bp]current page.north west) {#4};%
}
\newcommand{\pdffillrect}[4]{%
  \fill ([xshift=#1bp,yshift=-#2bp]current page.north west)
        rectangle ([xshift=#3bp,yshift=-#4bp]current page.north west);%
}
\newcommand{\pdfglyph}[5]{%
  \node[anchor=base west,inner sep=0pt,outer sep=0pt] at
    ([xshift=#1bp,yshift=-#2bp]current page.north west)
    {{#3\fontsize{#4bp}{#4bp}\selectfont #5}};%
}

% 一页式正式封面。静态文本位置严格取自 PDF 文字 baseline / vector line 坐标。
\newcommand{\makeopeningcover}{%
  \clearpage
  \thispagestyle{empty}
  \begingroup
  \xeCJKsetup{PunctStyle=plain}
  \begin{tikzpicture}[remember picture,overlay]
    % ---------- 静态文字 ----------
    % 所有静态字符均按用户 PDF 的逐字符 origin 坐标放置，避免 CJK 标点压缩造成偏移。
    % AUTO-GENERATED FROM MEASURED PDF CHARACTER ORIGINS
% Static cover characters only; dynamic fields are overlaid separately.
% Coordinates are PDF bp from the top-left origin.

\pdfglyph{90.024}{91.700}{\CoverDengAll}{14}{附}
\pdfglyph{104.064}{91.700}{\CoverDengAll}{14}{件}
\pdfglyph{121.580}{91.700}{\CoverDengAll}{14}{1}
\pdfglyph{225.650}{156.500}{\CoverSong}{18}{西}
\pdfglyph{243.650}{156.500}{\CoverSong}{18}{北}
\pdfglyph{261.650}{156.500}{\CoverSong}{18}{农}
\pdfglyph{279.650}{156.500}{\CoverSong}{18}{林}
\pdfglyph{297.650}{156.500}{\CoverSong}{18}{科}
\pdfglyph{315.650}{156.500}{\CoverSong}{18}{技}
\pdfglyph{333.650}{156.500}{\CoverSong}{18}{大}
\pdfglyph{351.650}{156.500}{\CoverSong}{18}{学}
\pdfglyph{261.650}{218.900}{\CoverSong}{18}{开}
\pdfglyph{279.650}{218.900}{\CoverSong}{18}{题}
\pdfglyph{297.650}{218.900}{\CoverSong}{18}{报}
\pdfglyph{315.650}{218.900}{\CoverSong}{18}{告}
\pdfglyph{120.020}{306.770}{\CoverHei}{22}{题}
\pdfglyph{141.980}{306.770}{\CoverHei}{22}{目}
\pdfglyph{163.940}{306.770}{\CoverHei}{22}{（}
\pdfglyph{185.900}{306.770}{\CoverHei}{22}{居}
\pdfglyph{207.948}{306.770}{\CoverHei}{22}{中}
\pdfglyph{229.908}{306.770}{\CoverHei}{22}{）}
\pdfglyph{240.954}{306.770}{\CoverHei}{22}{：}
\pdfglyph{120.020}{337.970}{\CoverHei}{22}{英}
\pdfglyph{141.980}{337.970}{\CoverHei}{22}{文}
\pdfglyph{163.940}{337.970}{\CoverHei}{22}{题}
\pdfglyph{185.900}{337.970}{\CoverHei}{22}{目}
\pdfglyph{207.948}{337.970}{\CoverHei}{22}{（}
\pdfglyph{229.996}{337.970}{\CoverHei}{22}{居}
\pdfglyph{251.956}{337.970}{\CoverHei}{22}{中}
\pdfglyph{273.916}{337.970}{\CoverHei}{22}{）}
\pdfglyph{284.962}{337.970}{\CoverHei}{22}{：}
\pdfglyph{146.660}{438.650}{\CoverSong}{15}{学}
\pdfglyph{191.570}{438.650}{\CoverSong}{15}{院}
\pdfglyph{214.130}{438.650}{\CoverSong}{15}{（}
\pdfglyph{229.130}{438.650}{\CoverSong}{15}{系}
\pdfglyph{251.570}{438.650}{\CoverSong}{15}{、}
\pdfglyph{266.675}{438.650}{\CoverSong}{15}{所}
\pdfglyph{281.675}{438.650}{\CoverSong}{15}{）}
\pdfglyph{146.660}{469.870}{\CoverSong}{15}{一}
\pdfglyph{161.660}{469.870}{\CoverSong}{15}{级}
\pdfglyph{176.660}{469.870}{\CoverSong}{15}{学}
\pdfglyph{191.660}{469.870}{\CoverSong}{15}{科}
\pdfglyph{214.130}{469.870}{\CoverTNR}{15}{/}
\pdfglyph{225.770}{469.870}{\CoverSong}{15}{类}
\pdfglyph{240.770}{469.870}{\CoverSong}{15}{别}
\pdfglyph{255.875}{469.870}{\CoverSong}{15}{（}
\pdfglyph{270.875}{469.870}{\CoverSong}{15}{领}
\pdfglyph{285.875}{469.870}{\CoverSong}{15}{域}
\pdfglyph{300.875}{469.870}{\CoverSong}{15}{）}
\pdfglyph{146.660}{501.070}{\CoverSong}{15}{研}
\pdfglyph{161.660}{501.070}{\CoverSong}{15}{究}
\pdfglyph{176.660}{501.070}{\CoverSong}{15}{生}
\pdfglyph{199.130}{501.070}{\CoverSong}{15}{姓}
\pdfglyph{214.130}{501.070}{\CoverSong}{15}{名}
\pdfglyph{236.570}{501.070}{\CoverSong}{15}{及}
\pdfglyph{259.250}{501.070}{\CoverSong}{15}{学}
\pdfglyph{274.250}{501.070}{\CoverSong}{15}{号}
\pdfglyph{146.660}{532.270}{\CoverSong}{15}{指}
\pdfglyph{169.100}{532.270}{\CoverSong}{15}{导}
\pdfglyph{191.570}{532.270}{\CoverSong}{15}{教}
\pdfglyph{214.130}{532.270}{\CoverSong}{15}{师}
\pdfglyph{244.130}{532.270}{\CoverSong}{15}{姓}
\pdfglyph{266.690}{532.270}{\CoverSong}{15}{名}
\pdfglyph{146.660}{563.470}{\CoverSong}{15}{开}
\pdfglyph{161.660}{563.470}{\CoverSong}{15}{题}
\pdfglyph{176.660}{563.470}{\CoverSong}{15}{论}
\pdfglyph{191.660}{563.470}{\CoverSong}{15}{证}
\pdfglyph{206.660}{563.470}{\CoverSong}{15}{小}
\pdfglyph{221.660}{563.470}{\CoverSong}{15}{组}
\pdfglyph{236.660}{563.470}{\CoverSong}{15}{组}
\pdfglyph{251.660}{563.470}{\CoverSong}{15}{长}
\pdfglyph{266.660}{563.470}{\CoverSong}{15}{姓}
\pdfglyph{281.660}{563.470}{\CoverSong}{15}{名}
\pdfglyph{146.660}{594.670}{\CoverSong}{15}{开}
\pdfglyph{161.660}{594.670}{\CoverSong}{15}{题}
\pdfglyph{176.660}{594.670}{\CoverSong}{15}{论}
\pdfglyph{191.660}{594.670}{\CoverSong}{15}{证}
\pdfglyph{206.660}{594.670}{\CoverSong}{15}{小}
\pdfglyph{221.660}{594.670}{\CoverSong}{15}{组}
\pdfglyph{236.660}{594.670}{\CoverSong}{15}{成}
\pdfglyph{251.660}{594.670}{\CoverSong}{15}{员}
\pdfglyph{266.660}{594.670}{\CoverSong}{15}{姓}
\pdfglyph{281.660}{594.670}{\CoverSong}{15}{名}
\pdfglyph{146.660}{703.300}{\CoverSong}{15}{研}
\pdfglyph{161.660}{703.300}{\CoverSong}{15}{究}
\pdfglyph{176.660}{703.300}{\CoverSong}{15}{生}
\pdfglyph{191.660}{703.300}{\CoverSong}{15}{通}
\pdfglyph{206.660}{703.300}{\CoverSong}{15}{过}
\pdfglyph{221.660}{703.300}{\CoverSong}{15}{开}
\pdfglyph{236.660}{703.300}{\CoverSong}{15}{题}
\pdfglyph{251.690}{703.300}{\CoverSong}{15}{论}
\pdfglyph{266.690}{703.300}{\CoverSong}{15}{证}
\pdfglyph{281.690}{703.300}{\CoverSong}{15}{日}
\pdfglyph{296.690}{703.300}{\CoverSong}{15}{期}
\pdfglyph{311.690}{703.300}{\CoverSong}{15}{：}
\pdfglyph{356.590}{703.300}{\CoverSong}{15}{年}
\pdfglyph{401.590}{703.300}{\CoverSong}{15}{月}
\pdfglyph{446.620}{703.300}{\CoverSong}{15}{日}


    % ---------- 学位层次行（硕/博可选，整体重新居中） ----------
    % 原 Word 的“博士（硕士）”为二选一占位。选择后将“年级线 + 级 + 学位层次 + 研究生学位论文”
    % 作为一个整体保持在 A4 页面水平中心。
    \pdffillrect{180.620}{189.620}{234.764}{190.460}
    \pdfglyph{234.770}{187.700}{\CoverSong}{18}{级}
    \ifReportDoctor
      \pdfglyph{252.770}{187.700}{\CoverSong}{18}{博}
    \else
      \pdfglyph{252.770}{187.700}{\CoverSong}{18}{硕}
    \fi
    \pdfglyph{270.770}{187.700}{\CoverSong}{18}{士}
    \pdfglyph{288.770}{187.700}{\CoverSong}{18}{研}
    \pdfglyph{306.770}{187.700}{\CoverSong}{18}{究}
    \pdfglyph{324.770}{187.700}{\CoverSong}{18}{生}
    \pdfglyph{342.770}{187.700}{\CoverSong}{18}{学}
    \pdfglyph{360.770}{187.700}{\CoverSong}{18}{位}
    \pdfglyph{378.770}{187.700}{\CoverSong}{18}{论}
    \pdfglyph{396.770}{187.700}{\CoverSong}{18}{文}

    % ---------- PDF 中其余 7 条实线（按 vector rect 精确坐标） ----------
    % 学院、学科、姓名学号、导师、组长、成员两行
    \pdffillrect{296.690}{440.350}{476.740}{441.070}
    \pdffillrect{323.350}{471.550}{480.930}{472.270}
    \pdffillrect{296.690}{502.750}{476.860}{503.470}
    \pdffillrect{289.130}{533.950}{476.740}{534.670}
    \pdffillrect{304.130}{565.150}{476.740}{565.870}
    \pdffillrect{304.130}{596.350}{476.740}{597.070}
    \pdffillrect{304.130}{627.580}{476.740}{628.300}

    % ---------- 可填写内容 ----------
    % 年级：居中放在线上方
    \pdfbasecenter{207.692}{187.700}{}{{\CoverTNR\fontsize{18bp}{18bp}\selectfont \ReportGrade}}

    % 题目：原 Word 只给标签，不画横线。填写时在标签右侧区域居中。
    \node[anchor=base,inner sep=0pt,outer sep=0pt,text width=285bp,align=center]
      at ([xshift=414bp,yshift=-306.770bp]current page.north west)
      {{\CoverSong\fontsize{15bp}{18bp}\selectfont \ReportTitle}};
    \node[anchor=base,inner sep=0pt,outer sep=0pt,text width=240bp,align=center]
      at ([xshift=427bp,yshift=-337.970bp]current page.north west)
      {{\CoverTNR\fontsize{13bp}{15bp}\selectfont \ReportTitleEN}};

    % 信息填写值：以每条线的中心位置对齐
    \pdfbasecenter{386.715}{438.650}{}{{\CoverSong\fontsize{15bp}{15bp}\selectfont \ReportDepartment}}
    \pdfbasecenter{402.140}{469.870}{}{{\CoverSong\fontsize{15bp}{15bp}\selectfont \ReportDiscipline}}
    \pdfbasecenter{386.775}{501.070}{}{{\CoverSong\fontsize{15bp}{15bp}\selectfont \ReportStudent\quad{\CoverTNR\ReportStudentID}}}
    \pdfbasecenter{382.935}{532.270}{}{{\CoverSong\fontsize{15bp}{15bp}\selectfont \ReportSupervisor}}
    \pdfbasecenter{390.435}{563.470}{}{{\CoverSong\fontsize{15bp}{15bp}\selectfont \ReportGroupLeader}}
    \pdfbasecenter{390.435}{594.670}{}{{\CoverSong\fontsize{15bp}{15bp}\selectfont \ReportGroupMembers}}

    % 日期填写值：中心点来自 PDF 中“年/月/日”前的空白区
    \pdfbasecenter{341.640}{703.300}{}{{\CoverTNR\fontsize{15bp}{15bp}\selectfont \ReportDefenseYear}}
    \pdfbasecenter{386.590}{703.300}{}{{\CoverTNR\fontsize{15bp}{15bp}\selectfont \ReportDefenseMonth}}
    \pdfbasecenter{431.605}{703.300}{}{{\CoverTNR\fontsize{15bp}{15bp}\selectfont \ReportDefenseDay}}
  \end{tikzpicture}
  \endgroup
  \null\newpage
}

% 新 PDF 已将开题论证小组成员与日期合并到同一封面页，因此不再生成第二页。
\newcommand{\makeopeningapproval}{}
